% ---------------------------------------------------------------------------
% macros.tex -- one macro per notation-bearing symbol.
%
% Why: renaming a symbol paper-wide becomes a one-line edit, and you can test
% whether a letter is still free by commenting out its macro and recompiling.
% Keep the notation table below in sync with the definitions.
%
%  symbol   macro     meaning                    type / domain
%  ---------------------------------------------------------------------
%  R        \R        real numbers               field
%  G        \Gr       communication graph        undirected, n nodes
%  L        \Lap      graph Laplacian of \Gr     R^{n x n}, PSD
%  lambda_2 \lamtwo   algebraic connectivity     >= 0
% ---------------------------------------------------------------------------

\newcommand{\R}{\mathbb{R}}
\newcommand{\N}{\mathbb{N}}
\newcommand{\Prob}{\mathbb{P}}
\newcommand{\Exp}{\mathbb{E}}

\newcommand{\Gr}{\mathcal{G}}
\newcommand{\Lap}{L}
\newcommand{\lamtwo}{\lambda_{2}}

\newcommand{\eps}{\varepsilon}          % never \epsilon: reads as \in
\newcommand{\emptyS}{\varnothing}       % never \emptyset: reads as 0
\newcommand{\norm}[1]{\left\lVert #1 \right\rVert}
\newcommand{\abs}[1]{\left\lvert #1 \right\rvert}
\newcommand{\ip}[2]{\langle #1, #2 \rangle}

\DeclareMathOperator{\tr}{tr}
\DeclareMathOperator{\rank}{rank}
\DeclareMathOperator*{\argmin}{arg\,min}
\DeclareMathOperator*{\argmax}{arg\,max}

\newcommand{\todo}[1]{\textbf{[TODO: #1]}}   % must be empty before submission
